% =====================================================================
%  Chapter 10: Time Series and Panel Data
%  Source: TimeSeries.tex (Overleaf stub, expanded) + new scaffold
%  NOTE: Section 10.1-10.4 need full development
% =====================================================================

\chapter{Time Series and Panel Data}
\label{chap_TimeSeries}

\section{Introduction and Key Concepts}
\label{sec_TSIntro}

Time series analysis studies the properties of sequences of observations
indexed by time. The key distinction from cross-sectional analysis is that
observations are not independent---the value at time $t$ is typically
correlated with past values (autocorrelation) and the statistical properties
of the series may change over time (non-stationarity).

\noindent\textbf{Why time series matters for regional analysis.}
Economic time series at the metropolitan and county level---employment,
income, housing prices, unemployment rates---are the primary inputs for
regional monitoring, forecasting, and policy evaluation. Understanding their
time-series properties is a prerequisite for valid statistical inference.

\subsection{Stationarity}

A time series $\{y_t\}$ is \textbf{weakly stationary}\index{stationarity}
if its mean, variance, and autocovariance structure do not change over time:
\begin{align}
  \E[y_t] &= \mu \quad\text{(constant mean)} \\
  \Var[y_t] &= \sigma^2 \quad\text{(constant variance)} \\
  \Cov[y_t, y_{t-k}] &= \gamma_k \quad\text{(depends only on lag }k\text{, not }t\text{)}
\end{align}

\subsection{Autocorrelation}

The \textbf{autocorrelation function} (ACF)\index{autocorrelation function}
at lag $k$ is:
\begin{equation}
  \rho_k = \frac{\Cov[y_t, y_{t-k}]}{\Var[y_t]} = \frac{\gamma_k}{\gamma_0}
  \;\in\; [-1,+1]
\end{equation}

The \textbf{partial autocorrelation function} (PACF) measures the
correlation at lag $k$ after controlling for all shorter lags.


\section{AR and ARMA Processes}
\label{sec_ARMA}

\subsection{The AR(1) Model}

The simplest time series model:
\begin{equation}
  y_t = \phi y_{t-1} + e_t, \qquad e_t \sim \text{iid}(0,\sigma^2)
\end{equation}

If $|\phi| < 1$: \emph{stationary}---the series mean-reverts.
If $\phi = 1$: \emph{unit root}---the series is a random walk and is
non-stationary.

\subsection{AR($p$) and ARMA($p$,$q$) Models}

The AR($p$) generalisation:
\begin{equation}
  y_t = \phi_1 y_{t-1} + \phi_2 y_{t-2} + \cdots + \phi_p y_{t-p} + e_t
\end{equation}

Adding $q$ moving-average terms gives the ARMA($p$,$q$) model:
\begin{equation}
  y_t = \sum_{i=1}^{p}\phi_i y_{t-i} + e_t + \sum_{j=1}^{q}\theta_j e_{t-j}
\end{equation}

The Box-Jenkins methodology selects $p$ and $q$ using the ACF and PACF
plots:
\begin{itemize}
  \item ACF cuts off at lag $q$, PACF decays: suggests MA($q$)
  \item PACF cuts off at lag $p$, ACF decays: suggests AR($p$)
  \item Both decay gradually: suggests ARMA($p$,$q$)
\end{itemize}

% TODO: Add R code for ARIMA modelling of regional employment series


\section{Non-Stationarity and Unit Roots}
\label{sec_UnitRoot}

\subsection{Spurious Regression}

If two time series are independently generated random walks (unit root
processes), regressing one on the other will typically produce a
statistically significant coefficient with high $R^2$---even though there
is no causal relationship. This is \textbf{spurious
regression}\index{spurious regression} \citep{Granger1974}.

\subsection{Augmented Dickey-Fuller Test}

The \textbf{Augmented Dickey-Fuller (ADF) test}\index{Dickey-Fuller test}
tests the null hypothesis of a unit root:
\begin{equation}
  \Delta y_t = \alpha + \beta t + \delta y_{t-1}
               + \sum_{j=1}^{k}\gamma_j \Delta y_{t-j} + e_t
\end{equation}
$H_0$: $\delta = 0$ (unit root); $H_1$: $\delta < 0$ (stationary).

\subsection{Cointegration and Error Correction}

Two non-stationary series are \textbf{cointegrated}\index{cointegration} if
a linear combination of them is stationary. Cointegration implies a
long-run equilibrium relationship; deviations from it are corrected over time
via an error correction mechanism:
\begin{equation}
  \Delta y_t = \alpha(y_{t-1} - \beta x_{t-1}) + \Gamma\Delta x_t + e_t
\end{equation}

% TODO: Add cointegration example with regional house prices and income


\section{Forecasting}
\label{sec_Forecasting}

\subsection{In-Sample Fit vs.\ Out-of-Sample Accuracy}

\begin{itemize}
  \item \textbf{In-sample fit} ($R^2$, AIC, BIC): measures how well the
    model fits the estimation sample. Subject to overfitting.
  \item \textbf{Out-of-sample accuracy} (RMSE, MAE on holdout data):
    measures genuine predictive performance.
\end{itemize}

\subsection{Forecast Evaluation}

The \textbf{Diebold-Mariano test} \citep{Diebold1995} compares the
predictive accuracy of two forecasting models:
\begin{equation}
  DM = \frac{\bar{d}}{\hat\sigma_{\bar{d}}/\sqrt{T}} \sim \mathcal{N}(0,1)
\end{equation}
where $\bar{d}$ is the mean loss differential between models.

% TODO: Expand principles of forecasting, including FRED regional data


\section{Panel Data Methods}
\label{sec_PanelData_TS}

Panel data methods were introduced in Section~\ref{sec_Panel} in the context
of causal identification. This section extends the treatment to dynamic
panels and the connection to time-series properties.

\subsection{Dynamic Panels}

When the lagged dependent variable $y_{i,t-1}$ is included as a regressor,
standard fixed effects estimation is inconsistent because $y_{i,t-1}$ is
correlated with $\alpha_i$ in the demeaned equation. The
\textbf{Arellano-Bond estimator}\index{Arellano-Bond estimator}
\citep{Arellano1991} uses lagged levels of $y$ as instruments for
$\Delta y_{i,t-1}$ in a first-differenced GMM framework.

\subsection{Spatial Panels}

Combining cross-sectional spatial dependence with time series produces
\textbf{spatial panel models}\index{spatial panel models}:
\begin{equation}
  y_{it} = \rho \sum_j w_{ij}y_{jt} + X_{it}'\beta + \alpha_i + \lambda_t + e_{it}
\end{equation}

This model allows simultaneous control for: (i) spatial spillovers via the
SAR term; (ii) unit-level unobserved heterogeneity via $\alpha_i$; and (iii)
common time shocks via $\lambda_t$.

% TODO: Add worked example with FRED metropolitan employment data

\renewcommand\bibname{Further reading}
\begin{subappendices}
\bibliographystyle{econometrica}
\bibliography{Literature/OverLord}
\IfFileExists{Literature/TimeSeries.bbl}{\chapter{Time Series Analysis}\label{chap_TimeSeries}
\section{Introduction and key concepts}
\section{Autocorrelation and non-stationarity}
\section{Principles of forecasting}



\renewcommand\bibname{Further reading}
\begin{subappendices}
\bibliographystyle{econometrica}
\bibliography{Literature/OverLord}
\IfFileExists{Literature/TimeSeries.bbl}{\chapter{Time Series Analysis}\label{chap_TimeSeries}
\section{Introduction and key concepts}
\section{Autocorrelation and non-stationarity}
\section{Principles of forecasting}



\renewcommand\bibname{Further reading}
\begin{subappendices}
\bibliographystyle{econometrica}
\bibliography{Literature/OverLord}
\IfFileExists{Literature/TimeSeries.bbl}{\chapter{Time Series Analysis}\label{chap_TimeSeries}
\section{Introduction and key concepts}
\section{Autocorrelation and non-stationarity}
\section{Principles of forecasting}



\renewcommand\bibname{Further reading}
\begin{subappendices}
\bibliographystyle{econometrica}
\bibliography{Literature/OverLord}
\IfFileExists{Literature/TimeSeries.bbl}{\input{Literature/TimeSeries.bbl}}{\typeout{Need to run bibtex}}
\end{subappendices} }{\typeout{Need to run bibtex}}
\end{subappendices} }{\typeout{Need to run bibtex}}
\end{subappendices} }{\typeout{Need bibtex on TimeSeries}}
\end{subappendices}
